\chapter{Conclusiones y recomendaciones}
\label{chap:conclusions_recommendations}

\section{Conclusiones generales}
\label{sec:general_conclusions}
% TODO: Sintetizar las principales conclusiones del proyecto, destacando los aportes técnicos y sociales logrados, y evaluando el cumplimiento del objetivo general de desarrollar un prototipo móvil para lectura automática de medidores en JAAPS rurales.

\section{Lecciones aprendidas del proceso de desarrollo}
\label{sec:lessons_learned}
% TODO: Reflexionar sobre el proceso de desarrollo completo: desafíos enfrentados durante la construcción del dataset, optimización de modelos para dispositivos móviles, integración de componentes, y experiencias de validación en campo con usuarios reales.

\section{Limitaciones técnicas y metodológicas}
\label{sec:limitations}
% TODO: Identificar y analizar las limitaciones del sistema desarrollado: restricciones del modelo OCR, dependencias de condiciones de captura, limitaciones de hardware móvil, alcance del dataset, y aspectos metodológicos que pudieron afectar los resultados.

\section{Recomendaciones para futuras investigaciones}
\label{sec:future_recommendations}
% TODO: Proponer líneas de investigación futuras: mejoras en arquitecturas de modelos, técnicas de few-shot learning para nuevos tipos de medidores, optimizaciones de hardware específico, y estudios longitudinales de adopción tecnológica en comunidades rurales.

\section{Posibilidades de escalabilidad y transferencia tecnológica}
\label{sec:scalability_transfer}
% TODO: Analizar el potencial de escalamiento del sistema: aplicabilidad a otros tipos de medidores, adaptación a diferentes contextos geográficos, consideraciones para implementación a gran escala, y estrategias de transferencia tecnológica a organizaciones de desarrollo rural.